% ------------------------------------------------------------------
%  Capitolo 22 — Studenti con DSA e altri bisogni speciali
%  Parte V
%  Ricerca di riferimento: ricerca 04 §8–11; ricerca 10 (da fare)
%  Voci bibliografiche principali: iss2022, tressoldi2001, wood2018, gersten2009, swanson1999, melbylervag2016, kuster2018, knouse2016, dupaul2012
%  Epigrafe: ✔ verificata sul testo (vedi ricerca 03 e registro, ricerca 00)
%  Scheda del capitolo: ricerca/06_indice-ragionato.md, capitolo 22
% ------------------------------------------------------------------
\chapter{Studenti con DSA e altri bisogni speciali}
\label{cap:dsa}

\epigrafe[Tre cose sono necessarie a chi studia: l'indole, l'esercizio, la disciplina.]{Tria sunt studentibus necessaria: natura, exercitium, disciplina.}{Ugo di San Vittore, Didascalicon III, 6}

\iniziale{Q}{uesto capitolo} \segnaposto{apertura: da scrivere — vedi la scheda nell'indice ragionato}.

\section{Che cosa sono i DSA, e chi li scopre all'università}

\segnaposto{da scrivere}

\section{Leggere lentamente in una lingua trasparente}

\segnaposto{da scrivere}

\section{Che cosa funziona e che cosa no}

\segnaposto{da scrivere}

\section{Compensare l'accesso, non il pensiero}

\segnaposto{da scrivere}

\section{Servizi di ateneo, ADHD e altri bisogni}

\segnaposto{da scrivere}

\begin{daricordare}
\segnaposto{il principio del capitolo in due righe}
\end{daricordare}

\begin{domande}
  \item \segnaposto{domanda di recupero su questo capitolo}
  \item \segnaposto{domanda di applicazione a una materia che studi}
  \item \segnaposto{domanda di ripresa su un capitolo precedente}
\end{domande}

\begin{sintesi}
  \item \segnaposto{punto di sintesi}
\end{sintesi}

\finecapitolo
